\documentclass[10pt]{article}
\usepackage[
  paperwidth=8.5in,
  paperheight=11in,
  left=2.45in,
  right=0.7in,
  top=0.75in,
  bottom=0.75in,
  headheight=14pt,
  headsep=18pt
]{geometry}

\usepackage{fontspec}
\usepackage{amsmath,amssymb}
\usepackage{graphicx}
\usepackage{xcolor}
\usepackage{enumitem}
\usepackage{fancyhdr}
\usepackage{tikz}

\definecolor{stewartred}{RGB}{190, 30, 35}

% Định nghĩa nhãn số phương trình hộp vuông đỏ
\newcommand{\eqtag}[1]{%
  \tikz[baseline={([yshift=-0.6ex]X.center)}]{%
    \node[fill=stewartred, text=white, font=\bfseries\footnotesize, inner xsep=4pt, inner ysep=3pt, rounded corners=1.2pt] (X) {#1};%
  }%
}

\pagestyle{fancy}
\fancyhf{}
\fancyheadoffset[L]{1.75in}
\fancyhead[L]{\small\textbf{292}\hspace{2.5em}\textbf{\textsf{CHƯƠNG 4}}\hspace{1.5em}\textsf{Các ứng dụng của đạo hàm}}
\renewcommand{\headrulewidth}{0pt}

\fancyfootoffset[L]{1.75in}
\fancyfoot[L]{%
  \parbox{\dimexpr\textwidth+1.75in\relax}{%
    \centering\fontsize{5.5pt}{7pt}\selectfont\color{black!60}%
    Copyright 2021 Cengage Learning. All Rights Reserved. May not be copied, scanned, or duplicated, in whole or in part. WCN 02-200-203\\
    Do quyền kỹ thuật số, một số nội dung bên thứ ba có thể bị ẩn khỏi sách điện tử và/hoặc từng chương. Đánh giá biên tập khẳng định rằng nội dung bị ẩn không ảnh hưởng đáng kể đến trải nghiệm học tập tổng thể.%
  }%
}

\begin{document}

\noindent
Trước khi chứng minh định lý này, chúng ta có thể thấy rằng nó hoàn toàn hợp lý khi diễn giải theo mặt hình học. Hình 3 và 4 biểu diễn các điểm $A(a, f(a))$ và $B(b, f(b))$ trên đồ thị của hai hàm số khả vi. Hệ số góc của cát tuyến $AB$ là

\vspace{3pt}
\noindent\makebox[\linewidth]{\eqtag{3}\hfill $\displaystyle m_{AB} = \frac{f(b) - f(a)}{b - a}$\hfill\phantom{\eqtag{3}}}
\vspace{5pt}

\noindent
đây chính là biểu thức ở vế phải của Phương trình 1. Vì $f'(c)$ là hệ số góc của tiếp tuyến tại điểm $(c, f(c))$, nên Định lý Giá trị Trung bình, dưới dạng được cho bởi Phương trình 1, khẳng định rằng có ít nhất một điểm $P(c, f(c))$ trên đồ thị mà tại đó hệ số góc của tiếp tuyến bằng hệ số góc của cát tuyến $AB$. Nói cách khác, tồn tại một điểm $P$ mà tại đó tiếp tuyến song song với cát tuyến $AB$. (Hãy hình dung một đường thẳng ở xa luôn song song với $AB$ và tịnh tiến dần về phía $AB$ cho đến khi chạm vào đồ thị lần đầu tiên.)

\vspace{10pt}
\noindent
\begin{minipage}[t]{0.48\textwidth}
  \centering
  \includegraphics[width=\linewidth]{images/p0327_fig1.png}\\[5pt]
  {\small\bfseries\textsf{HÌNH 3}}
\end{minipage}\hfill
\begin{minipage}[t]{0.48\textwidth}
  \centering
  \includegraphics[width=\linewidth]{images/p0327_fig2.png}\\[5pt]
  {\small\bfseries\textsf{HÌNH 4}}
\end{minipage}
\vspace{12pt}

\noindent%
\begin{tikzpicture}[remember picture, overlay]
  \node[anchor=north east, inner sep=0pt, xshift=-16pt, yshift=4pt] at (0,0) {%
    \begin{minipage}{140pt}
      \centering
      \includegraphics[width=\linewidth]{images/p0327_fig3.png}\\[6pt]
      {\small\bfseries\textsf{HÌNH 5}}
    \end{minipage}%
  };
\end{tikzpicture}%
\textbf{\textsf{\textcolor{stewartred}{CHỨNG MINH}}}\quad Ta áp dụng Định lý Rolle cho một hàm số mới $h$ được xác định là hiệu giữa $f$ và hàm số có đồ thị là đường cát tuyến $AB$. Sử dụng Phương trình 3 và phương trình đường thẳng theo điểm và hệ số góc, ta thấy phương trình của đường thẳng $AB$ có thể viết dưới dạng

\vspace{3pt}
\[
y - f(a) = \frac{f(b) - f(a)}{b - a}(x - a)
\]
\vspace{-4pt}
\noindent\makebox[\linewidth]{\text{hoặc là}\hfill $\displaystyle y = f(a) + \frac{f(b) - f(a)}{b - a}(x - a)$\hfill\phantom{\text{hoặc là}}}
\vspace{5pt}

\noindent
Vì vậy, như minh họa trong Hình 5,

\vspace{3pt}
\noindent\makebox[\linewidth]{\eqtag{4}\hfill $\displaystyle h(x) = f(x) - f(a) - \frac{f(b) - f(a)}{b - a}(x - a)$\hfill\phantom{\eqtag{4}}}
\vspace{5pt}

\noindent
Trước tiên, ta phải kiểm tra xem $h$ có thỏa mãn ba giả thiết của Định lý Rolle hay không.

\begin{enumerate}[label=\textbf{\arabic*.}, leftmargin=1.5em, itemsep=5pt, topsep=4pt]
  \item Hàm số $h$ liên tục trên $[a, b]$ vì nó là tổng của $f$ và một đa thức bậc nhất, cả hai hàm này đều liên tục.
  \item Hàm số $h$ khả vi trên $(a, b)$ vì cả $f$ và đa thức bậc nhất đều khả vi. Trên thực tế, ta có thể tính đạo hàm $h'$ trực tiếp từ Phương trình 4:
  \[
  h'(x) = f'(x) - \frac{f(b) - f(a)}{b - a}
  \]
  (Lưu ý rằng $f(a)$ và $[f(b) - f(a)]/(b - a)$ là các hằng số.)
\end{enumerate}

\end{document}